% ============================================
% PAPER OUTLINE - ĐỀ XUẤT SƠ BỘ BÀI BÁO
% ============================================
% A Multi-Layer Machine Learning Framework for DDoS Detection in IoT Networks
% with Adaptive Decision Engine

\documentclass[conference]{IEEEtran}
\usepackage{cite}
\usepackage{amsmath,amssymb,amsfonts}
\usepackage{algorithmic}
\usepackage{graphicx}
\usepackage{textcomp}
\usepackage{xcolor}

\begin{document}

\title{A Multi-Layer Machine Learning Framework for DDoS Detection in IoT Networks with Adaptive Decision Engine\\
{\large \textit{Paper Outline - Đề xuất sơ bộ}}}

\author{\IEEEauthorblockN{Nguyen Tat Thang, Tran Dang Cong}
\IEEEauthorblockA{Faculty of Information Technology, Dai Nam University, Ha Noi\\
Email: congtd@dainam.edu.vn}}

\maketitle

\section{ABSTRACT}
The rapid proliferation of Internet of Things (IoT) devices has significantly increased the attack surface of modern network infrastructures, making Distributed Denial-of-Service (DDoS) attacks a critical security concern. This paper proposes a multi-layer machine learning framework for intelligent DDoS detection in IoT networks, integrating protocol-level analysis and network traffic analytics with an adaptive decision engine. The proposed framework performs intrusion detection on both MQTT communication traffic and general network flows using benchmark datasets, including the MQTT IoT Intrusion Detection Dataset and UNSW-NB15. To address severe class imbalance and improve model robustness, data preprocessing techniques such as feature encoding, normalization, Synthetic Minority Over-sampling Technique (SMOTE), and random under-sampling are employed. Multiple machine learning algorithms are evaluated, where LightGBM and XGBoost demonstrate superior performance for multi-class and binary detection tasks, respectively. Furthermore, an adaptive decision engine incorporating probability calibration, threat-level classification, and dynamic threshold optimization is developed to enhance real-time detection reliability. Experimental results show that the proposed framework achieves F1-scores of 93.06\% for MQTT traffic and 95.23\% for network traffic, outperforming baseline models. The system also demonstrates low inference latency suitable for real-time deployment. These findings highlight the effectiveness of combining multi-source traffic learning with adaptive decision strategies to strengthen IoT network security.

\textbf{Keywords:} DDoS detection, Multi-layer machine learning, IoT network security

\section{I. INTRODUCTION}

\subsection{Background and Motivation}
\begin{itemize}
    \item IoT devices proliferation and security challenges
    \item DDoS attacks as critical threat to IoT networks
    \item Limitations of traditional detection methods
    \item Need for intelligent, adaptive detection systems
\end{itemize}

\subsection{Research Gap}
\begin{itemize}
    \item Existing works focus on single protocol or single dataset
    \item Limited integration of protocol-level (MQTT) and network-level detection
    \item Lack of adaptive decision mechanisms for real-time deployment
    \item Class imbalance problem not adequately addressed
\end{itemize}

\subsection{Research Objectives}
\begin{itemize}
    \item Develop multi-layer ML framework for both MQTT and network traffic
    \item Address class imbalance through advanced preprocessing
    \item Design adaptive decision engine with probability calibration
    \item Evaluate performance on benchmark datasets
    \item Achieve real-time detection capability
\end{itemize}

\subsection{Contributions}
\begin{itemize}
    \item Novel multi-layer framework integrating MQTT and network traffic analysis
    \item Comprehensive preprocessing pipeline with SMOTE and under-sampling
    \item Adaptive decision engine with Isotonic Calibration and Attack Bias Logic
    \item Extensive evaluation showing F1-scores of 93.06\% (MQTT) and 95.23\% (Network)
    \item Low-latency system suitable for real-time IoT deployment
\end{itemize}

\subsection{Paper Organization}
Sections: Introduction → Literature Review → Methodology → Results → Conclusion

\section{II. LITERATURE REVIEW / THEORETICAL FRAMEWORK}

\subsection{DDoS Attacks in IoT Networks}
\begin{itemize}
    \item Types of DDoS attacks (volumetric, protocol, application layer)
    \item IoT-specific vulnerabilities
    \item MQTT protocol security challenges
    \item Impact on network infrastructure
\end{itemize}

\subsection{Machine Learning for Intrusion Detection}
\begin{itemize}
    \item Traditional ML approaches (Decision Trees, Random Forest, SVM)
    \item Ensemble methods (XGBoost, LightGBM, CatBoost)
    \item Deep learning approaches (CNN, LSTM, Autoencoders)
    \item Comparison of algorithms for network intrusion detection
\end{itemize}

\subsection{Class Imbalance Handling}
\begin{itemize}
    \item SMOTE (Synthetic Minority Over-sampling Technique)
    \item Random under-sampling
    \item Hybrid approaches
    \item Impact on model performance
\end{itemize}

\subsection{Probability Calibration}
\begin{itemize}
    \item Isotonic Regression
    \item Platt Scaling
    \item Calibration for reliable predictions
    \item Applications in security systems
\end{itemize}

\subsection{Related Works}
\begin{itemize}
    \item MQTT-based intrusion detection systems
    \item Network traffic analysis for DDoS detection
    \item Hybrid detection frameworks
    \item Limitations of existing approaches
\end{itemize}

\section{III. RESEARCH METHODOLOGY / RESEARCH MODEL}

\subsection{Overall Framework Architecture}
\begin{itemize}
    \item Four-phase pipeline: Preprocessing → Training → Optimization → Decision Engine
    \item Multi-layer approach: MQTT layer + Network layer
    \item Integration strategy
\end{itemize}

\subsection{Phase 1: Data Preprocessing}

\subsubsection{MQTT Traffic Preprocessing}
\begin{itemize}
    \item Dataset: MQTT IoT Intrusion Detection Dataset
    \item Features: 42 features including protocol-specific metrics
    \item Missing value handling
    \item Feature encoding (Label Encoding for categorical features)
    \item Feature scaling (StandardScaler)
    \item Train-test split (70-30)
    \item Class distribution analysis
\end{itemize}

\subsubsection{Network Traffic Preprocessing}
\begin{itemize}
    \item Dataset: UNSW-NB15
    \item Features: 42 network flow features
    \item Categorical feature encoding
    \item Numerical feature normalization
    \item Train-test split
    \item Attack type distribution
\end{itemize}

\subsection{Phase 2: Model Training and Comparison}

\subsubsection{Class Imbalance Handling}
\begin{itemize}
    \item SMOTE for minority class over-sampling
    \item Random under-sampling for majority class
    \item Balanced dataset creation
    \item Impact on model performance
\end{itemize}

\subsubsection{Machine Learning Algorithms}
\begin{itemize}
    \item Random Forest: Ensemble of decision trees
    \item XGBoost: Gradient boosting with regularization
    \item LightGBM: Histogram-based gradient boosting
    \item CatBoost: Categorical feature handling
\end{itemize}

\subsubsection{Training Strategy}
\begin{itemize}
    \item Cross-validation approach
    \item Hyperparameter settings
    \item Training time analysis
    \item Model comparison metrics
\end{itemize}

\subsection{Phase 3: Model Optimization}

\subsubsection{Hyperparameter Tuning}
\begin{itemize}
    \item RandomizedSearchCV methodology
    \item Parameter search space
    \item Optimization objectives
    \item Best parameter selection
\end{itemize}

\subsubsection{Model Validation}
\begin{itemize}
    \item Testing on raw (unbalanced) dataset
    \item Generalization capability assessment
    \item Overfitting prevention
\end{itemize}

\subsection{Phase 4: Adaptive Decision Engine}

\subsubsection{Probability Calibration}
\begin{itemize}
    \item Isotonic Regression calibration
    \item Calibration set (20\% of test data)
    \item Evaluation set (80\% of test data)
    \item Reliability diagram analysis
\end{itemize}

\subsubsection{Attack Bias Logic}
\begin{itemize}
    \item Confidence threshold mechanism
    \item Attack class prioritization
    \item Decision rules:
        \begin{itemize}
            \item If confidence $\geq$ threshold: Select highest probability class
            \item If confidence $<$ threshold: Check attack classes
            \item If attack probability $>$ 0.25: Select attack class
            \item Otherwise: Select highest probability class
        \end{itemize}
\end{itemize}

\subsubsection{Threat Level Classification}
\begin{itemize}
    \item ROC curve analysis
    \item Optimal threshold determination
    \item Three-level classification: NORMAL, MEDIUM, CRITICAL
    \item Threshold values: 0.4 (MQTT), 0.1064 (Network)
\end{itemize}

\subsection{Evaluation Metrics}
\begin{itemize}
    \item Accuracy, Precision, Recall, F1-score
    \item Confusion matrix analysis
    \item ROC-AUC score
    \item Processing time and latency
\end{itemize}

\section{IV. RESULTS AND DISCUSSION}

\subsection{Phase 1: Data Analysis Results}
\begin{itemize}
    \item MQTT dataset characteristics
    \item Network dataset characteristics
    \item Class distribution visualization
    \item Feature importance analysis
\end{itemize}

\subsection{Phase 2: Model Comparison Results}

\subsubsection{MQTT Traffic Detection}
\begin{itemize}
    \item Performance comparison table (4 algorithms)
    \item LightGBM achieves best F1-score: 93.06\%
    \item Confusion matrix analysis
    \item Per-class performance metrics
\end{itemize}

\subsubsection{Network Traffic Detection}
\begin{itemize}
    \item Performance comparison table
    \item XGBoost achieves best F1-score: 95.23\%
    \item Binary classification results
    \item Attack detection rate analysis
\end{itemize}

\subsection{Phase 3: Optimization Results}
\begin{itemize}
    \item Hyperparameter tuning improvements
    \item Performance on raw dataset
    \item Generalization capability
    \item Training vs. testing performance
\end{itemize}

\subsection{Phase 4: Decision Engine Results}

\subsubsection{Calibration Effectiveness}
\begin{itemize}
    \item Before vs. after calibration comparison
    \item Reliability improvement
    \item Calibration curve analysis
\end{itemize}

\subsubsection{Attack Bias Logic Impact}
\begin{itemize}
    \item False negative reduction
    \item Attack detection improvement
    \item Trade-off analysis
\end{itemize}

\subsubsection{Threat Level Classification}
\begin{itemize}
    \item Threshold optimization results
    \item Three-level classification accuracy
    \item Real-time decision making capability
\end{itemize}

\subsection{System Performance}
\begin{itemize}
    \item Inference latency: $<$ 10ms per sample
    \item Throughput analysis
    \item Resource utilization
    \item Scalability assessment
\end{itemize}

\subsection{Comparison with Baseline}
\begin{itemize}
    \item Comparison with existing methods
    \item Performance improvements
    \item Advantages and limitations
\end{itemize}

\subsection{Discussion}
\begin{itemize}
    \item Key findings interpretation
    \item Practical implications for IoT security
    \item Limitations and challenges
    \item Future research directions
\end{itemize}

\section{V. RECOMMENDATIONS AND CONCLUSION}

\subsection{Recommendations}
\begin{itemize}
    \item Deployment strategies for IoT networks
    \item Integration with existing security infrastructure
    \item Continuous learning and model updating
    \item Threshold tuning for specific environments
\end{itemize}

\subsection{Conclusion}
\begin{itemize}
    \item Summary of contributions
    \item Key achievements:
        \begin{itemize}
            \item Multi-layer framework for MQTT and network traffic
            \item F1-scores: 93.06\% (MQTT), 95.23\% (Network)
            \item Adaptive decision engine with calibration
            \item Real-time detection capability
        \end{itemize}
    \item Significance for IoT security
    \item Future work directions
\end{itemize}

\section{VI. REFERENCES}
\begin{itemize}
    \item IoT security and DDoS attack papers
    \item Machine learning for intrusion detection
    \item SMOTE and class imbalance handling
    \item Probability calibration methods
    \item MQTT protocol security
    \item UNSW-NB15 dataset papers
    \item XGBoost, LightGBM, CatBoost papers
\end{itemize}

\section{VII. APPENDICES (if any)}
\begin{itemize}
    \item Detailed algorithm pseudocode
    \item Additional experimental results
    \item Hyperparameter search space details
    \item Extended performance tables
\end{itemize}

\vspace{1cm}
\noindent\textbf{Note:} This outline provides the structure for the full paper. Each section will be expanded with detailed content, equations, figures, and tables in the complete version.

\end{document}
